\section{Motivation}
\label{sec:motivation}

Modern key--value stores --- RocksDB, LevelDB, Cassandra, HBase --- are built
on the \emph{log-structured merge-tree} (LSM-tree)~\citep{oneil1996lsm}.
The LSM-tree is a write-optimised data structure: incoming writes accumulate in an in-memory
buffer and are flushed to disk as immutable Sorted String Tables (SSTables)
at the top level, while a background compaction process merges and promotes
files to successively larger levels.
The structure converts random writes into
sequential I/O at the cost of \emph{read amplification}: a single lookup
may have to consult an SSTable at every level.

To avoid touching each SSTable on disk, every file carries a small in-memory
\emph{filter} in its metadata.
The filter answers a membership-style question
with one-sided error: false negatives are forbidden, and the false positive
rate is bounded by~$\varepsilon$. If the filter says ``empty'', the engine
skips the file with zero disk I/O. The Bloom filter~\citep{bloom1970space}
fills this role for point lookups and is the default in production
LSM-trees~\citep{monkey2017,dong2021rocksdb}. Bloom filters, however, answer
a range query only by probing every key in the interval --- time linear in
the range length, and impractical at LSM scale~\citep{zhang2018surf}.
Range-shaped workloads
--- short scans, prefix lookups over hierarchical keys (file paths, S2 cells,
URLs), and time-window queries --- form a non-trivial share of
production LSM traffic~\citep{cao2020facebook}.

For this regime we need a \emph{range filter}: given a query interval
$[a, b]$ of length at most~$\mathcal{L}$, decide whether $[a, b] \cap S = \emptyset$
under the same one-sided error guarantee.
The information-theoretic limit of this problem is settled.
\citet{goswami2015approximate} prove a worst-case lower bound of
$\log_2(\mathcal{L}/\varepsilon)$ bits per key and give a matching
construction --- a locality-preserving hash followed by an
\emph{Exact Range Emptiness} (ERE) backend --- but their construction
is, in their successors' own words, ``mainly theoretic in nature and
thus very complicated to implement''~\citep{grafite2024}.

The bound is worst-case, however, and \citet{vaidya2022snarf} make
the escape clause explicit: it holds only \emph{without further
assumptions on the data or workload}, so a filter tuned to a specific
input regime can do better. Existing practical filters take one of
two paths along this clause: either accept an FPR that holds only on
uncorrelated, in-distribution queries (SuRF, SNARF), or pay for a
distribution-free guarantee with heavier per-query
work~\citep{luo2020rosetta,grafite2024}. Real LSM keys offer a third
path that none of them exploits: production workloads cluster into
dense blocks separated by sparse gaps~\citep{cao2020facebook}. This
thesis pursues that third path. First, we tighten the constants in
the Goswami-style ERE backend, narrowing the practical gap to the
worst-case lower bound. Second, where the data distribution permits,
we detect dense clusters and handle them exactly, going below the
bound on benign inputs while preserving the worst-case guarantee on
adversarial ones.
